\chapter*{KATA PENGANTAR}
\addcontentsline{toc}{chapter}{KATA PENGANTAR}

Puji syukur kami panjatkan ke hadirat Tuhan Yang Maha Esa karena atas rahmat dan karunia-Nya, saya dapat menyelesaikan penulisan Laporan Tugas Akhir yang berjudul "Pengembangan Sistem Manajemen Karyawan Terpusat Berbasis Integrasi Perangkat Kontrol Akses pada Bangunan Gedung Cerdas". Laporan ini disusun sebagai salah satu syarat untuk menyelesaikan program Sarjana di Sekolah Teknik Elektro dan Informatika, Institut Teknologi Bandung.

Pada kesempatan ini, saya ingin menyampaikan terima kasih kepada: 
\begin{enumerate}
    \item Dr. Fadhil Hidayat, S.Kom., M.T., selaku pembimbing saya, yang telah memberikan bimbingan, dukungan, dan masukan yang sangat berharga selama proses penyusunan laporan ini.
    \item Bapak I Gusti Bagus Baskara Nugraha, S.T., M.T., PhD. sebagai ketua program studi Sistem dan Teknologi Informasi Institut Teknologi Bandung.
    \item Muhammad Rifa Ansyari dan Natanael Steven Simangunsong sebagai rekan tim tugas akhir
    \item Keluarga saya, yang selalu memberikan dukungan moral dan motivasi tanpa henti selama saya menempuh pendidikan ini.
    \item Kepada Kepala Seksi Data DKST, pengelola operasional Gedung IIP, serta seluruh staf yang telah memberikan izin wawancara, akses data, dan dukungan fasilitas pengujian gerbang.
    \item Teman-teman dan rekan-rekan di Institut Teknologi Bandung, yang telah memberikan bantuan, diskusi, dan semangat selama perjalanan akademik saya.
    \item Para pengguna dan penguji sistem yang telah memberikan feedback konstruktif untuk perbaikan sistem ini.
    \item Semua pihak yang telah membantu, baik secara langsung maupun tidak langsung, dalam penyelesaian tugas akhir ini.
\end{enumerate}
\begin{flushright}
    Bandung, 29 Mei 2026\\
    
    Penulis \\
    \vspace{1cm}
    Axelius Davin
\end{flushright}
